\documentclass[10pt,a4paper]{article}
\usepackage[margin=22mm]{geometry}
\usepackage{amsmath,amssymb,booktabs,array}
\usepackage[colorlinks=true,urlcolor=blue]{hyperref}
\usepackage{microtype}
\setlength{\parindent}{0pt}
\setlength{\parskip}{5pt}
\title{RM9: Approved section-speed observer\\Current-task history audit and its limits}
\author{Method-first ramp-metering review}
\date{5 October 2026}
\begin{document}
\maketitle
\textbf{Disposition.} Observer C is approved; UG4-O is closed. This bounded R2
pass establishes conditional information sufficiency, not a positive
equal-history/useful-action certificate or universal impossibility. R3 admits
no rescue assumption. R5 screens alternatives only; UG5 requires a user
disposition before a changed task is formulated. No traffic experiment is run.

\section{Approved observation, unchanged traffic law}
For a mainline cell of length $L$, write lane masses $n_1,n_2$ in vehicle
equivalents, total $N=n_1+n_2$, density $k_\ell=n_\ell/L$ and storage $M_\ell$.
Public FD parameters are $v_\ell,w_\ell,C_\ell$ and
$k^{\rm jam}_\ell=M_\ell/L$. Define $c_\ell=v_\ell/L$,
$\gamma_\ell=w_\ell/L$ and the constitutive equilibrium rate
\begin{equation}
g_\ell(n)=\min\{c_\ell n,C_\ell,\gamma_\ell(M_\ell-n)\}.
\end{equation}
With $U^{\rm eq}_\ell(k)=f^{\rm eq}_\ell(k)/k$ for occupied lanes, C observes
\begin{equation}
V=\frac{n_1U^{\rm eq}_1(k_1)+n_2U^{\rm eq}_2(k_2)}{N},\qquad
\frac{NV}{L}=g_1(n_1)+g_2(n_2),\quad N>0.
\label{eq:c}
\end{equation}
Empty-lane contributions are zero without evaluating $0/0$; an empty cell's
speed is unavailable. This is a homogeneous representative-equilibrium closure,
not a detector or receiving-limited physical-speed theorem. The right side is
\emph{not} the actual accepted face rate, which is observed separately.

Preserve the adopted ordering: bounded lateral transfer conserving mass, then
recompute forward sending and receiving, then conservative forward/merge updates.
The Default positive-priority allocator redistributes unused receiving; the
mainline and ramp share their forward receiver. Meter permission, actual
release and actual downstream merge entry remain different where storage lies
between them. No new traffic state or law follows from the observer approval.

\textbf{TST qualification.} C describes traffic present in a section. It does
not compute time spent in the whole system. For time step $h$ and all-stream
present population $P^s$,
\begin{equation}
J_h=h\sum_{s=0}^{S-1}P^s.
\end{equation}
$P^s$ includes mainline, pre/post-meter ramp storage and external waiting, with
the same arrivals and clearance rule for every method. Normalization by a fixed
common total mass does not change action ranking. Do not replace this objective
by a cell length divided by $V$; the FD closure need not describe actual blocked
motion. The same-region time--space speed identity is a different statistic.
\newpage

\section{Branch-local information and global ambiguity}
With $z=n_1$ and observed $N$, feasible splits lie in
\begin{equation}
I_N=[\max(0,N-M_2),\min(M_1,N)],\qquad
G_N(z)=g_1(z)+g_2(N-z)=NV/L.
\end{equation}
Use F for the free branch, P for a capacity plateau and H for the congested
branch. Away from branch boundaries, the lane-rate slopes are respectively
$c_\ell,0,-\gamma_\ell$. The slope of $G_N$ is the difference of lane slopes:
\begin{center}
\begin{tabular}{lccc}\toprule
Lane 1 / lane 2 & F & P & H\\\midrule
F & $c_1-c_2$ & $c_1$ & $c_1+\gamma_2$\\
P & $-c_2$ & $0$ & $\gamma_2$\\
H & $-(\gamma_1+c_2)$ & $-\gamma_1$ & $\gamma_2-\gamma_1$\\
\bottomrule\end{tabular}
\end{center}
A nonzero slope determines at most one candidate \emph{inside that known
branch domain}. Equal free-flow mobilities or equal congested wave coefficients
make the respective same-branch measurement constant; plateau/plateau is also
constant when that branch actually exists. These are algebraic conditions,
not permission to impose equal parameters to manufacture a collision.

\textbf{Global warning.} Each $g_\ell$ is concave; $G_N$ is concave and
piecewise affine. Roots can occur on opposite sides of its maximum, even when
each involved branch has a nonzero slope. Flat intervals can give a continuum.
Thus a free-flow inversion is a global identification result only when all
other feasible branches are excluded by public state/domain information or
completed observations. We do not silently observe branch membership.
Nonconcave changes would require a new law, not use of this argument.

\textbf{Plateau warning.} A positive-width P branch requires
\begin{equation}
C_\ell<\frac{c_\ell\gamma_\ell M_\ell}{c_\ell+\gamma_\ell}.
\end{equation}
For the ordinary triangular FD, equality holds and P is only a peak point.
FB0030's triangular capacity relation does not justify a positive interval
with both lanes on P. A separately capped trapezoid could be a legitimate
public model if independently supported; it has not been adopted here.

\textbf{Actual faces.} Even a compatible section-speed root must reproduce
completed actual aggregate face counts/rates, ramp balances, all other cells
and prior controls under the adopted lateral/forward law. A capped face alone
does not remove the separate section-speed information. Future action-induced
flow cannot be used to decide which root existed before choosing that action.
\newpage

\section{The complete-history requirement}
Let $x^s$ denote the complete modeled mainline/ramp state, $p$ public parameters,
$b^s$ all declared forcing/boundary allocation, and $u^s$ prior meter permissions.
The deterministic adopted update is
\begin{equation}
x^{s+1}=F_B(x^s,b^s,u^s;p).
\end{equation}
Let $O_C$ include stored-state inventories/section speeds and completed labeled
face and ramp observations. At the decision epoch, define the feasible initial
state set using the \emph{whole common declared history} $H_*$:
\begin{equation}
\mathcal X(H_*)=\{x^0:\ x^{s+1}=F_B(x^s,b^s,u^s;p),\
O_C(x^s,\text{completed fluxes})=O^s_{\rm admitted}\ \text{for all }s\leq *\}.
\end{equation}
All candidates obey storage, geometry, conservation and public forcing. This
defines the proof obligation; it is not a numerical estimator or a solved
site-specific constraint system.

\textbf{Conditional sufficiency lemma.} If the admitted information uniquely
identifies the complete state at any common past epoch $a\leq *$, then all later
states up to the decision are uniquely determined. Proof: apply the public
deterministic update once to the known state and inputs; repeat by induction.
For any common future schedule and forcing, the same induction gives identical
future states, exit histories and TST, hence identical action rankings.

This includes a fully known complete initialization (including a genuinely
known empty system), or an earlier epoch where every state coordinate is
uniquely reconstructed. Queue formation, capacity caps and later snapshot
ambiguity do not invalidate this result. A uniquely invertible local branch
in one cell does not by itself identify the whole system.

\textbf{Unresolved domain.} If only aggregate initial information is known
and multiple complete states survive \emph{all} declared observations, the
lemma does not settle the task. Mixed/congested branches may permit snapshot
ambiguity, but no genuine site/initialization-supported pair has been
established here that survives the coupled full history and changes useful
common actions. Conversely, no global impossibility theorem has been proved
for every such domain. Do not make the history short or forcing private just
to obtain multiplicity.

If a noninjective lateral update maps two different pre-states to the same
\emph{complete} state before actions act, with all remaining inputs equal,
their lost coordinate cannot create a subsequent action distinction. A failure
of a memoryless total-density closure is also not sufficient: a stronger
aggregate-history method can reconstruct an otherwise missing coordinate.
\newpage

\section{Action and gradient status; next physical comparison}
The required positive claim would need two genuine feasible states with
identical $H_*$, distinguishable useful actions from the same admissible
open-loop family, and complete affected-population time accounting. None is
admitted by this pass. Therefore R4 has no new certificate to propagate into
a claim of useful gradients. Existing branch-wise conceptual derivatives are
conditional; no finite-difference, cross-fidelity or one-step usefulness test
has been conducted in this methodology review.

R3 declines unsupported symmetry, invented plateau widths, weakened/noisy
speed observations, omitted waiting, private forcing and extra lanes introduced
only to defeat an inverse. These could be legitimate features of another
source-grounded task, but need explicit scientific admission and cannot be
silent repairs of this one.

R5's bounded source comparison supports physical questions, not positivity:
\begin{itemize}
\item Laval, Cassidy and Daganzo's inspected merge model represents bounded
acceleration and associated voids, even at a site without an acceleration lane.
Its ramp-queue-dependent entry speed is not automatically hidden when that
queue is known. No exact RM9 history collision follows.
\item Srivastava's 2016 dissertation, Chapter 3.3.1, changes congested demand
and adds a jam-demand parameter while retaining supply and density
conservation. This is an aggregate-law competitor, in a queue-discharge task;
the separate lane-drop model in Chapter 4 requires additional merge modeling
for a ramp extension. Neither is copied into RM9.
\item A real finite acceleration lane may require position/urgency and changed
geometry. Gap-informed metering changes sensing and temporal action semantics.
Both require their own task comparison, not an unchanged-law counterexample.
\end{itemize}
\textbf{Recommendation at UG5.} Accept the scoped absence of an established
positive result for this bounded current-task pass, and discuss a targeted
recovery/acceleration-family review beginning with the strongest credible
aggregate demand/history competitor. Add state only if necessary to represent
independently supported information that this competitor cannot reconstruct.
This is a proposed next scope, not an adopted law, successor or universal
rejection of the current family. The companion findings present fair alternatives.

\section*{Primary provenance and reviewer boundary}
FD source: \href{https://repository.tudelft.nl/file/File_2c4359e7-5e5a-4651-964f-6f3610ab89ae}{FB0030, final primary lane-model PDF}, capacity relation printed 1768.
Physical merge mechanism: \href{https://bpb-us-e1.wpmucdn.com/sites.gatech.edu/dist/f/3571/files/2023/03/merges.pdf}{Laval, Cassidy and Daganzo, inspected merge paper}, Sections 2/4.
Aggregate competitor: \href{https://escholarship.org/content/qt7dx478g4/qt7dx478g4.pdf}{Srivastava, UCI dissertation 2016}, selected Chapters 3/4.
Observer provenance remains in the R1 source matrix. Equations and history
lemmas above are reviewer deductions under the explicitly approved closure,
not literature theorems for the combined adapted task. Search/access and claim
ledgers retain version and coverage limits. No source establishes the combined
hidden-information, useful-action and useful-gradient chain.
\end{document}
